% BHU PYQ -- Professional Exam Template
\documentclass[11pt,a4paper]{article}

\usepackage[a4paper,top=2.5cm,bottom=2.5cm,left=2.8cm,right=2.8cm,headheight=14pt]{geometry}
\usepackage{amsmath,amssymb}
\usepackage[T1]{fontenc}
\usepackage[utf8]{inputenc}
\usepackage{lmodern}
\usepackage{microtype}
\usepackage{fancyhdr}
\usepackage{enumitem}
\usepackage{hyperref}
\usepackage{lastpage}
\usepackage{parskip}

\hypersetup{colorlinks=true,urlcolor=black,linkcolor=black,
  pdftitle={B.Sc. (Hons.) Botany Sem-I BOB-101 -- BHU PYQ}}

\pagestyle{fancy}
\fancyhf{}
\renewcommand{\headrulewidth}{0.4pt}
\renewcommand{\footrulewidth}{0.4pt}
\fancyhead[L]{\small   \textit{Banaras Hindu University}}
\fancyhead[R]{\small   \textit{Botany Paper: BOB-101}}
\fancyfoot[C]{\small Page  \thepage\ of \pageref{LastPage}}


\newlist{parts}{enumerate}{2}
\setlist[parts,1]{label=(\alph*),leftmargin=*,itemsep=4pt,topsep=3pt}
\setlist[parts,2]{label=(\roman*),leftmargin=*,itemsep=2pt,topsep=2pt}

\newcommand{\pts}[1]{\hfill   \textit{[#1]}}


\begin{document}


\begin{center}
  {\large\bfseries BANARAS HINDU UNIVERSITY}\\[2pt]
  {
\normalsize B.Sc. (Hons.) Semester-I Examination, 2022-23}\\[6pt]
  \rule{0.8\linewidth}{0.5pt}\\[6pt]
  {\large\bfseries Botany}\\[4pt]
  {\large\bfseries Paper: BOB-101 --- Cryptogams}\\[4pt]
  {
\normalsize Time: Three Hours \hfill Maximum Marks: 70}
\end{center}

\rule{\linewidth}{0.4pt}

\smallskip



\noindent   \textit{Instructions: Answer five questions, selecting at least two each from Sections A and B and Section-C which is compulsory. The figures in the right hand margin indicate marks.}

\bigskip


\begin{center}
   \textbf{SECTION A}
\end{center}
\medskip

\medskip


\noindent   \textbf{Question 1.} Describe the characteristic features of Xanthophyceae. Give an illustrated account of reproduction in   \textit{Vaucheria}.\pts{15}

\medskip


\noindent   \textbf{Question 2.} Write short notes on any two of the following:\pts{$7.5  \times 2 = 15$}

\begin{parts}
  \item Sexual reproduction in   \textit{Ectocarpus}
  \item Morphology and reproduction of   \textit{Scytonema}
  \item Asexual reproduction in   \textit{Volvox}
\end{parts}

\medskip


\noindent   \textbf{Question 3.} With the help of a labeled diagram, give a brief account of the morphology and sexual reproductive structure of   \textit{Marchantia}.\pts{15}

\medskip


\noindent   \textbf{Question 4.} Write short notes on any two of the following:\pts{$7.5  \times 2 = 15$}

\begin{parts}
  \item Sexual reproduction in   \textit{Anthoceros}
  \item Cap cell and zoospore formation in   \textit{Oedogonium}
  \item Sporophyte of   \textit{Funaria}
\end{parts}

\bigskip

\begin{center}
   \textbf{SECTION B}
\end{center}
\medskip

\medskip


\noindent   \textbf{Question 5.} Give an illustrated account of various types of spore formation in   \textit{Puccinia}.\pts{15}

\medskip


\noindent   \textbf{Question 6.} Write short notes on any two of the following:\pts{$7.5  \times 2 = 15$}

\begin{parts}
  \item Reproduction in   \textit{Rhizopus}
  \item General characteristics of Basidiomycotina
  \item Morphology of basidiocarp of   \textit{Agaricus}
\end{parts}

\medskip


\noindent   \textbf{Question 7.} With the help of a suitable diagram, give a brief account of the morphology and reproduction in   \textit{Selaginella}.\pts{15}

\medskip


\noindent   \textbf{Question 8.} Write short notes on any two of the following:\pts{$7.5  \times 2 = 15$}

\begin{parts}
  \item T.S. stem of   \textit{Equisetum} through internode
  \item   \textit{Cercospora}
  \item T.S. rachis of   \textit{Pteris}
\end{parts}

\bigskip

\begin{center}
   \textbf{SECTION C}
\end{center}
\medskip

\medskip


\noindent   \textbf{Question 9.} Answer the following (not more than 50 words each):\pts{$1  \times 10 = 10$}

\begin{parts}
  \item What is the difference between elater and pseudoelater?
  \item What is the function of vallecular canal?
  \item Name the fungal group which possesses Crozier.
  \item What is the function of receptacle in   \textit{Sargassum}?
  \item What are the xerophytic characteristics of   \textit{Equisetum} stem?
  \item Name the Zygospore-forming fungi.
  \item What is dikaryotic phase?
  \item Name the pigments present in Cyanophyceae members.
  \item Name the reserved food material present in the Phaeophyceae.
  \item What is heterokont flagella?
\end{parts}

\vfill
\rule{\linewidth}{0.4pt}

\begin{center}\small
\href{https://bhu-exam.vercel.app/}{Click Here}\\[4pt]\href{https://me-aryan.vercel.app/}{Click Here}\\[4pt]\href{https://quashan-me.vercel.app/}{Click Here}
\end{center}

\end{document}
